\clearpage
\section{21 September 2026: Applying reviewed parent and land decisions}

The production pipeline now applies the supported allocations from the
preceding ten-parent review. Bedford Square's four contemporaneous filings form
one 877-unit economic parent, and Starhill's two filings form one 570-unit
parent. Dupont retains its single 381-unit filing with corrected land.
Membership and parcel allocations are applied together. Administrative housing
counts and filing dates remain the source measurements.

\begin{center}
\small
\begin{tabular}{lrrr}
\hline
Development & Units & Ground (sq ft) & Earlier floor area (sq ft) \\
\hline
Dupont & 381 & 20,901 & 0 \\
Bedford Square & 877 & 177,520.79 & 14,946 \\
Starhill & 570 & 70,063 & 80,370 \\
\hline
\end{tabular}
\end{center}

The correction also locates existing buildings within the included ground.
At Bedford, recorded west housing parcels total 116,870.79 square feet; the
complete earlier east lot contributes 60,650 once. The full 175,875 square feet
of the former Sears building stays on the excluded retained parcel, as shown
by current PLUTO and supported by the DEC records and deed boundaries. The
housing parent retains only the earlier east auto-center's 14,946 square feet,
so its existing floor-area ratio is 0.08419. The earlier residential FAR is
2.43. Dupont's recorded ground combines 9,011 and 11,890 square feet, compared
with 12,194 previously; both reference parcels have zero existing building area.

At Starhill, the declaration excludes retained lots 129 and 162. The included
strip from old lot 162 measures 44.5 by 100 feet, or 4,450 square feet. The
remaining 65,613 square feet of the adopted 70,063 comes from old lot 134.
The parcel subdivision assigns all 33,750 square feet of the old lot-162
building to retained land. The development therefore retains old lot 134's
80,370 square feet, giving existing FAR 1.14711. Applying the 19v2 residential
FARs of 6.02 and 3.44 to their respective included areas gives 5.85613.

The historical canonical panel contains 1,814 parents, four fewer, with the
same 129,335 units. The post-policy panel remains 909 parents and 61,319 units.
The weighting sample has 557 historical and 312 post-policy parents.
Comparisons preserve every constituent's units, original and refiling dates,
and sample inclusion. Every other parent's existing panel values and the full
post-policy site-characteristics dataset are unchanged.

The same calibration specification produces positive weights from 0.1400 to
1.6830, summing to 312. Effective sample size is 492.9, compared with 495.4
before the correction; the largest standardized balance error is below
$6 \times 10^{-12}$. These are refreshed inputs to the existing pilot, whose
specification remains fixed while the broader land review continues.

The refreshed audit has 129 flagged parents out of 869, compared with 135 out
of 873. The six former Bedford/Starhill rows become two supported parents;
Dupont had already passed its documentary review. Bedford's source explicitly
resolves DOF transaction 99999, a November 2023 correction to lot 46's length
covered by the inspected successor map and later recorded deed description.
The raw flag remains available; other physical transactions are not waived.
All other parents retain their preceding pass/flag status.

Recorded boundaries support these observed development grounds. Starhill's
total uses a later administrative area measurement of that boundary; it is
not a February 2020 survey. Later records do not establish ownership at initial
filing or a legal wage-assessment unit. Eagle/West, 355 Exterior, and 1580 Story
were open at this stage. The following entry records the additional original
evidence and production corrections that resolve these three allocations.

\textit{Reproduction.} Root \texttt{make}, \texttt{make dof-site-review}, and
\texttt{make logbook}. The September 21 implementation note links the original
sources and verification.
